% ═══════════════════════════════════════════════════════════════════
%  Webographie
% ═══════════════════════════════════════════════════════════════════
\pageLiminaire{Webographie}

Les ressources ci-dessous ont été consultées durant la réalisation du projet.
Les dates indiquées correspondent à la dernière consultation.

\vspace{0.4cm}

\begin{enumerate}[leftmargin=1.1cm,itemsep=6pt]

  \item \textbf{NestJS} --- documentation officielle du cadriciel de l'API
  (modules, gardes, intercepteurs, filtres d'exceptions, limitation de débit).
  \\ \url{https://docs.nestjs.com/}

  \item \textbf{Prisma} --- documentation de l'ORM : schéma, migrations,
  transactions, modes de suppression référentielle.
  \\ \url{https://www.prisma.io/docs}

  \item \textbf{PostgreSQL} --- documentation des fonctions de verrouillage
  consultatif, primitive au c\oe ur du moteur de réservation.
  \\ \url{https://www.postgresql.org/docs/current/explicit-locking.html}

  \item \textbf{React} --- documentation officielle de la bibliothèque
  d'interface du back-office.
  \\ \url{https://react.dev/}

  \item \textbf{React Native} --- documentation officielle du cadriciel de
  l'application mobile.
  \\ \url{https://reactnative.dev/docs/getting-started}

  \item \textbf{Expo} --- documentation de la chaîne d'outils mobile : stockage
  sécurisé, polices, images, compilation native.
  \\ \url{https://docs.expo.dev/}

  \item \textbf{TypeScript} --- manuel de référence du langage.
  \\ \url{https://www.typescriptlang.org/docs/}

  \item \textbf{Vite} --- documentation de l'outil de construction du
  back-office.
  \\ \url{https://vite.dev/guide/}

  \item \textbf{OWASP} --- \textit{Password Storage Cheat Sheet} et
  \textit{Authentication Cheat Sheet} : recommandations sur le hachage des mots
  de passe et la défense contre l'énumération de comptes.
  \\ \url{https://cheatsheetseries.owasp.org/}

  \item \textbf{JSON Web Tokens} --- introduction au format et à ses usages.
  \\ \url{https://jwt.io/introduction}

  \item \textbf{Argon2} --- spécification de la fonction de hachage retenue
  (RFC 9106).
  \\ \url{https://datatracker.ietf.org/doc/html/rfc9106}

  \item \textbf{IANA Time Zone Database} --- base des fuseaux horaires,
  référence de la zone \texttt{Africa/Casablanca}.
  \\ \url{https://www.iana.org/time-zones}

  \item \textbf{date-fns} --- documentation de la bibliothèque de manipulation
  de dates et de son extension pour les fuseaux horaires.
  \\ \url{https://date-fns.org/docs/Getting-Started}

  \item \textbf{OpenAPI / Swagger} --- spécification du format de documentation
  d'API et de son interface interactive.
  \\ \url{https://swagger.io/specification/}

  \item \textbf{Docker} --- documentation de la conteneurisation et de Docker
  Compose.
  \\ \url{https://docs.docker.com/}

  \item \textbf{GitHub Actions} --- documentation de la chaîne d'intégration
  continue.
  \\ \url{https://docs.github.com/en/actions}

  \item \textbf{Jest} --- documentation du cadriciel de tests.
  \\ \url{https://jestjs.io/docs/getting-started}

  \item \textbf{Scrum Guide} --- guide de référence de la méthode SCRUM.
  \\ \url{https://scrumguides.org/}

  \item \textbf{UML} --- spécification du langage de modélisation unifié
  publiée par l'OMG.
  \\ \url{https://www.omg.org/spec/UML/}

  \item \textbf{Lucidchart} --- plate-forme de création des diagrammes UML du
  présent rapport.
  \\ \url{https://www.lucidchart.com/}

  \item \textbf{CNDP} --- Commission Nationale de contrôle de la protection des
  Données à caractère Personnel : loi 09-08 et droits des personnes concernées.
  \\ \url{https://www.cndp.ma/}

  \item \textbf{i18next} --- documentation de la bibliothèque
  d'internationalisation employée côté mobile.
  \\ \url{https://www.i18next.com/}

\end{enumerate}
